%----------------------------------------------------------------------------------------
% LOGARITMISCHE FUNCTIES, VERGELIJKINGEN EN ONGELIJKHEDEN. Alle opgaven zijn eigen werk.
%
% Opbouw: Start, leerstof (theorie, werkwijze of uitleg, voorbeeld, korte oefening), en op een
% nieuwe pagina de Opdrachten (door elkaar, stijgende moeilijkheid), Koppel en Test jezelf.
% Labels van dit hoofdstuk beginnen met h3-.
%----------------------------------------------------------------------------------------

\mychapter{3}{Logaritmische functies}

\lkr{Leerplan B+S: LPD 7-12 (exponentiële en logaritmische functies), het deel over de logaritmische functie, logaritmische vergelijkingen en ongelijkheden. De inverse functie komt uit de cursus Algemene functies. Het grootste struikelblok zijn de bestaansvoorwaarden: laat ze altijd eerst opschrijven, en de oplossingen op het einde ertegen controleren.}

In hoofdstuk 2 was de logaritme een getal. Nu bekijken we $x \mapsto \log_a x$ als functie. Ze is de inverse van de exponentiële functie: ze draait $a^x$ terug. Daarmee los je ook vergelijkingen en ongelijkheden op waarin de onbekende in een logaritme staat.

\begin{oefeningen}[Start]
\oef{h3-start-1}

\oef{h3-verken}
\end{oefeningen}

%========================================================================================
\section{De logaritmische functie}

\begin{theorie}
Voor $a > 0$ met $a \neq 1$ is $\formule{f(x) = \log_a x}$ de \textbf{\textcolor{accent}{logaritmische functie}} met grondtal $a$. Ze is de inverse functie van $a^x$, met
\[\formule{\dom f = \left]0, +\infty\right[} \qquad \formule{\ber f = \R}.\]
Haar grafiek gaat door $(1, 0)$ en $(a, 1)$, en heeft de $y$-as als \textbf{\textcolor{accent}{verticale asymptoot}}. Ze is \textbf{stijgend} als $a > 1$ en \textbf{dalend} als $0 < a < 1$.
\end{theorie}

\naast{%
De grafiek van $\log_2 x$ is het spiegelbeeld van die van $2^x$ om de rechte $y = x$. Domein en bereik wisselen: $2^x$ heeft alleen positieve beelden, dus $\log_2 x$ bestaat alleen voor positieve $x$.

Een logaritmische functie stijgt heel traag: $\log_2 x$ bereikt $10$ pas bij $x = 1024$. Dicht bij $0$ daalt ze onbeperkt: $\log_2 0{,}001 \approx -10$.}{%
\begin{tikzpicture}[x=.5cm,y=.5cm,font=\scriptsize]
  \rooster{-3}{1}{5}{-3}{1}{5}
  \draw[grijs,dashed] (-3,-3) -- (5.2,5.2) node[above,text=black] {$y = x$};
  \kromme[dashed]{-3}{2.32}{2^(\x)}
  \kromme{.13}{5}{ln(\x)/ln(2)}
  \vol{(1,0)} \vol{(2,1)}
  \node[right,text=accent] at (4.2,1.6) {$\log_2 x$};
  \node[left,text=accent] at (1.9,4.6) {$2^x$};
\end{tikzpicture}}

\voorbeeld{}{Schets de grafiek van $g(x) = \log_2(x - 3) + 1$. Geef het domein, de asymptoot en de nulwaarde.\label{vb:h3-grafiek}}{%
\naast{%
\uitleg{De grafiek van $\log_2 x$ schuift 3 naar rechts en 1 naar boven; de asymptoot schuift mee naar rechts.}
\begin{stappenvrij}{4.6cm}{3cm}
(1, 0) \mapsto (4, 1) & Het vaste punt. \\
\dom g = \left]3, +\infty\right[ & $x - 3 > 0$. \\
\text{asymptoot } x = 3 & \\
\log_2(x - 3) = -1 & Nulwaarde. \\
x - 3 = \tfrac12 \iff x = \tfrac72 & Definitie.
\end{stappenvrij}}{%
\begin{tikzpicture}[x=.5cm,y=.5cm,font=\scriptsize]
  \rooster{-1}{1}{7}{-3}{1}{3}
  \draw[grijs,dashed] (3,-3) -- (3,3.3);
  \kromme{3.07}{7}{ln(\x-3)/ln(2)+1}
  \vol{(4,1)} \vol{(3.5,0)}
\end{tikzpicture}}}

\begin{oefeningen}
\oef{h3-k-schets}
\end{oefeningen}

\voorbeeld{}{Bepaal het domein van $f(x) = \ln(4 - x^2)$.\label{vb:h3-domein}}{%
\uitleg{Van een logaritme moet het argument strikt positief zijn.}
\begin{stappenbreed}
4 - x^2 > 0 \iff x^2 < 4 & Bestaansvoorwaarde. \\
\dom f = \left]-2, 2\right[ & Tussen de nulwaarden van de bergparabool.
\end{stappenbreed}}

\begin{oefeningen}
\oef{h3-grafiek}

\oef{h3-k-domein}
\end{oefeningen}

%========================================================================================
\section{Logaritmische vergelijkingen}

Een logaritmische functie is overal stijgend of overal dalend, dus injectief: $\log_a A = \log_a B \iff A = B$, op voorwaarde dat $A$ en $B$ positief zijn.

\voorbeeld{}{Los op: $\log_3(2x - 1) = 2$.\label{vb:h3-vergelijking}}{%
\uitleg{Schrijf eerst de bestaansvoorwaarde op. Gebruik dan de definitie: de logaritme is de exponent.}
\begin{stappenbreed}
2x - 1 > 0 \iff x > \tfrac12 & Bestaansvoorwaarde. \\
2x - 1 = 3^2 = 9 & Definitie van de logaritme. \\
x = 5 & Voldoet aan $x > \frac12$.
\end{stappenbreed}}

\begin{oefeningen}
\oef{h3-k-def}
\end{oefeningen}

\voorbeeld{}{Los op: $\log_2(x + 1) + \log_2(x - 1) = 3$.\label{vb:h3-twee}}{%
\uitleg{Breng het linkerlid samen tot één logaritme met de productregel. Dat mag alleen waar beide logaritmen bestaan: daarom eerst de bestaansvoorwaarden.}
\begin{stappenbreed}
x + 1 > 0 \text{ en } x - 1 > 0 \iff x > 1 & Bestaansvoorwaarden. \\
\log_2\bigl((x + 1)(x - 1)\bigr) = 3 & Productregel. \\
x^2 - 1 = 8 \iff x = 3 \text{ of } x = -3 & Definitie. \\
x = 3 & $-3$ voldoet niet aan $x > 1$.
\end{stappenbreed}}

\begin{werkwijze}
Schrijf de \textbf{bestaansvoorwaarden} op: elk argument van een logaritme moet strikt positief zijn. Breng de vergelijking met de rekenregels in de vorm $\log_a A = c$ (dan is $A = a^c$) of $\log_a A = \log_a B$ (dan is $A = B$). Los op en houd alleen de oplossingen die aan de bestaansvoorwaarden voldoen.
\end{werkwijze}

\begin{oefeningen}
\oef{h3-vergelijking}

\oef{h3-k-twee}
\end{oefeningen}

%========================================================================================
\section{Ongelijkheden}

\begin{theorie}
\begin{stelling}{Ongelijkheden met machten en logaritmen}
Als $\formule{a > 1}$: \qquad $\formule{a^p < a^q \iff p < q}$ \qquad en \qquad $\formule{\log_a x < \log_a y \iff x < y}$.

Als $\formule{0 < a < 1}$, keert het ongelijkheidsteken om. Bij de logaritmen zijn $x, y > 0$.
\end{stelling}
\end{theorie}

Dat volgt uit het stijgen of dalen van de functies. Ook als je beide leden van een ongelijkheid door een negatief getal deelt, zoals $\ln 0{,}8$, keert het teken om.

\voorbeeld{}{Een medicijn wordt per uur met $20$\,\% afgebroken. Vanaf wanneer is er minder dan $10$\,\% van over?\label{vb:h3-ongelijkheid}}{%
\uitleg{Na $t$ uur is er nog een fractie $0{,}8^t$ over. Neem $\ln$ van beide leden; $\ln$ stijgt, dus het teken blijft. Daarna deel je door $\ln 0{,}8 < 0$: het teken keert om.}
\begin{stappenbreed}
0{,}8^t < 0{,}1 & \\
t \ln 0{,}8 < \ln 0{,}1 & $\ln$ van beide leden. \\
t > \frac{\ln 0{,}1}{\ln 0{,}8} \approx 10{,}3 & Delen door een negatief getal.
\end{stappenbreed}
Na ongeveer $10{,}3$ uur is er minder dan $10$\,\% over.}

\begin{oefeningen}
\oef{h3-k-expong}
\end{oefeningen}

\voorbeeld{}{Los op: $\log_3(x - 2) < 2$.\label{vb:h3-logongelijkheid}}{%
\uitleg{Bestaansvoorwaarde eerst. Schrijf dan $2$ als logaritme met grondtal 3; omdat $3 > 1$ blijft het teken.}
\begin{stappenbreed}
x - 2 > 0 \iff x > 2 & Bestaansvoorwaarde. \\
\log_3(x - 2) < \log_3 9 \iff x - 2 < 9 & $2 = \log_3 9$, grondtal $> 1$. \\
x \in \left]2, 11\right[ & Beide voorwaarden samen.
\end{stappenbreed}}

\begin{oefeningen}
\oef{h3-ongelijkheid}

\oef{h3-k-logong}
\end{oefeningen}

%========================================================================================
\opdrachtenkop

\begin{oefeningen}
\needspace{13\baselineskip}
\oef{h3-o1}

\oef{h3-o11}

\needspace{14\baselineskip}
\oef{h3-o2}

\needspace{12\baselineskip}
\oef{h3-o12}

\oef{h3-o13}

\oef{h3-o3}

\needspace{10\baselineskip}
\oef{h3-o14}

\needspace{18\baselineskip}
\oef{h3-o15}

\oef{h3-o16}

\oef{h3-o4}

\oef{h3-o19}

\oef{h3-o7}

\needspace{14\baselineskip}
\oef{h3-o6}

\oef{h3-o17}

\oef{h3-o20}

\oef{h3-o5}

\oef{h3-o22}

\oef{h3-o8}

\oef{h3-o18}

\oef{h3-o21}

\oef{h3-o23}

\oef{h3-o10}

\oef{h3-o24}

\oef{h3-o26}

\oef{h3-o9}

\oef{h3-o25}
\end{oefeningen}

%----------------------------------------------------------------------------------------
% KOPPEL: verbanden met andere leerstof
%----------------------------------------------------------------------------------------
\needspace{20\baselineskip}
\opdrachtsectie{Koppel}

\lkr{Algemene functies, Rijen en Integralen: kijk na wat al gezien is. Het pakket legt de volgorde niet vast.}

\begin{oefeningen}
\oef{h3-koppel-1}

\oef{h3-koppel-2}

\oef{h3-koppel-3}
\end{oefeningen}

%----------------------------------------------------------------------------------------
\opdrachtsectie{Test jezelf}

\begin{oefeningen}
\oef{h3-t1}

\oef{h3-t2}

\oef{h3-t3}

\oef{h3-t4}

\oef{h3-t5}

\oef{h3-t6}

\oef{h3-t7}
\end{oefeningen}

\herhaalschema{%
  \fv{oef:h3-t1}{vb:h3-grafiek}{oef:h3-grafiek}%
  \fv{oef:h3-t2}{vb:h3-domein}{oef:h3-k-domein}%
  \fv{oef:h3-t3}{vb:h3-vergelijking}{oef:h3-k-def}%
  \fv{oef:h3-t4}{vb:h3-twee}{oef:h3-k-twee}%
  \fv{oef:h3-t5}{vb:h3-ongelijkheid}{oef:h3-k-expong}%
  \fv{oef:h3-t6}{vb:h3-logongelijkheid}{oef:h3-k-logong}%
  \fv{oef:h3-t7}{vb:h3-grafiek}{oef:h3-o3}}

\ifzelfstudie
\section*{Maak je eigen samenvatting}
Sluit dit boek en schrijf in je schrift: domein, bereik, vaste punten en asymptoot van $\log_a x$ voor $a > 1$ en $0 < a < 1$, het verband met $a^x$, de werkwijze voor logaritmische vergelijkingen met de bestaansvoorwaarden, en wanneer het teken van een ongelijkheid omkeert. Vergelijk daarna met de kaders.
\fi
